% TOP_PROBLEM: {"id": 3387, "title": "Singmaster's conjecture", "category_id": 1, "category": {"id": 1, "name": "number_theory", "display_name": "Number Theory", "description": "Properties of integers, prime numbers, Diophantine equations.", "slug": "number-theory", "order_index": 1, "created_at": "2026-07-31T15:26:25.670Z"}, "status": "open", "problem_number": "OPG-60034", "set_id": 12}
% FIRST_POSTED: 2026-10-01
% ENTRY_KIND: statement_only
% UnsolvedMath ID 3387; source catalogue code OPG-60034
% !TeX program = lualatex
% Definition and sources only; this file is not an LLM solution attempt.
\documentclass[11pt,a4paper]{article}
\usepackage[margin=25mm]{geometry}
\usepackage{fontspec}
\setmainfont{lmroman10-regular.otf}[BoldFont=lmroman10-bold.otf,
 ItalicFont=lmroman10-italic.otf,BoldItalicFont=lmroman10-bolditalic.otf]
\usepackage{amsmath,amssymb,mathtools}
\usepackage{unicode-math}
\setmathfont{latinmodern-math.otf}
\AtBeginDocument{\let\setminus\smallsetminus}
\usepackage{xurl,hyperref}
\hypersetup{colorlinks=true,urlcolor=blue}
\setcounter{secnumdepth}{0}
\setlength{\parindent}{0pt}
\setlength{\parskip}{5pt}
\setlength{\emergencystretch}{3em}
\begin{document}
\section*{Singmaster's conjecture}\label{3387}
{\small UnsolvedMath ID 3387 · OPG-60034 · Number Theory · OpenGarden}
\subsection{Definitions and mathematical statement}
For $n\ge0$ and $0\le k\le n$, the binomial coefficient is
$\binom nk=n!/(k!(n-k)!)$. Pascal's triangle is counted by positions
$(n,k)$; in particular symmetric positions $(n,k)$ and $(n,n-k)$
count separately when they are different. For an integer $a>1$, put
\[
 N(a)=\#\{(n,k):n\ge0,\ 0\le k\le n,\ \binom nk=a\}.
\]
Singmaster's conjecture asserts the existence of an absolute finite
constant $C$ such that $N(a)\le C$ for every integer $a>1$.
The value one is excluded because it occurs along both unbounded
edges of Pascal's triangle.

Each individual $N(a)$ is finite: an occurrence with $a>1$ has
$1\le k\le n-1$ and $\binom nk\ge n$, so $n\le a$.
The conjecture asks for a bound independent of $a$, which is much
stronger than this individual finiteness. It does not specify a
particular numerical value of $C$ or require finding the smallest
one. A disproof must produce entries whose numbers of occurrences
are unbounded, rather than just a single entry occurring more often
than a proposed small bound. The convention about symmetric entries
fixes the multiplicity being measured.
\subsection{Short English statement}
Is there one finite bound on how many positions of Pascal’s triangle can contain any fixed integer greater than one?
\subsection{Sources}
\begin{itemize}
\item[S1] ULAM.AI and the UnsolvedMath Contributors, supplied problems.json, record 3387 (OPG-60034), OpenGarden collection. Catalogue identity and source transcription; CC BY 4.0.
\url{https://huggingface.co/datasets/ulamai/UnsolvedMath}.
\item[S2] Open Problem Garden, Singmaster's conjecture. Original problem and discussion; the mathematical formulation below is expanded from the supplied source record.
\url{http://www.openproblemgarden.org/op/singmasters_conjecture}.
\end{itemize}
\end{document}
